%!TEX root = ./../../thesis-trier.tex

Our convention for function composition deviates from the one adopted in permutation group theory
or in some approaches to the algebraic theory of automata. As this can be confusing at times, we feel the urge to justify and comment on that.
The notation and terminology introduced in the present section 
will not be included in the notation index or the index of terminology,
as it is only used here.

\IN{right action}
\IN{left action}
\IN{Leibniz order}
\IN{anti-Leibniz order} 
\IN{diagrammatic order}
When applying a function $f \colon X \to Y$ to some $x \in X$,
we can write the argument to the right, i.e., $f(x)$, and $f$ acts on the left (left action),
or we can write the argument to the left, i.e., $(x)f$ ($xf$ for short), and $f$ acts from the right (right action).
In the present work, we use the left-action convention (introduced on page~\pageref{def:function-application}).
Furthermore, when defining the composition of two functions
$f \colon X \to Y$ and $g \colon Y \to Z$, either the function on the right-hand side
can be applied first (right-most first, as defined here, see page~\pageref{def:fun_comp}), also
called the \emph{Leibniz order}, or the function on the left-hand side, i.e.,
define a composition operator by $(f ; g)(x) = g(f(x))$ for $x \in X$ (left-most first), also called the \emph{anti-Leibniz} or \emph{diagrammatic order}.

In~\cite{nlab:composition} also the notation $f^{\ast}g$ for $g \circ f$
was mentioned as pulling back~$g$ along~$f$ to the composite. Here,
the asterik operator makes $f \colon X \to Y$
into a function $f^{\ast} \colon Z^Y \to Z^X$
defined by $(f^{\ast}(g))(x) = g(f(x))$ for $x \in X$.
So, conceptually, rather than regarding composition as a binary operation
on functions, here we rather consider each function as an operator
acting via composition (this is reminiscent of letting each element of a monoid
act via multiplication from left or right, see~\cite[Prop.~4.25]{Pin2022}).
In the wild we have never seen anyone using this notation solely
for the purpose of having a notation for function composition,
but either in rather specialized contexts as operators on special
function spaces as, for example, in differential geometry~\cite[p.~77 \& p.~89]{spivak1971calculus},
as the composition operator in functional analysis~\cite{DAnielloMaiuriello2020_surveyCompositionOperators},
or in the context of Koopman Operator Theory~\cite{DograRedman2020_OptimizingNNsKoopman}.

Now, the following combinations of conventions are conceivable.

\begin{enumerate}
	
	\item Right-action and left-most first (also called ``postfix convention'').
	
    With these conventions, %$f$ applied to $x$ is often denoted by $xf$	so that
	we have $x(f ; g) = (xf)g$. This combination of conventions was
	in wide use in the early days of the theory of automata, 
	for example in the books~\cite{Eilenberg1974,Eilenberg1976,Ginzburg68,Pin86}. 
	It has the advantage of being
	in harmony with reading 
	a diagram $X \xrightarrow{f} Y \xrightarrow{g} Z$
	for the composition from left to right.
	In~\cite[p.~1]{Eilenberg1974},
	it is even noted that $x \in X$ can be interpreted as a function
	from some singleton set into~$X$ with value~$x$
	and then application of~$f$ to~$x$ can be written as function composition
	and modeled by diagrams depicting compositions. 
	However, these conventions seemed to be out of fashion in automata theory
	at least since the mid-nineties.
	In permutation group theory, this combination of conventions 
	is often used~\cite{cameron_1999,dixon1996permutation,wielandt1964finite}, even taken a step further by using
	the \emph{exponential notation} for function application, i.e., writing $x^f$
	for $f$ applied to $x$, see for example~\cite{Kurzweil04}.
		
	\item Left-action and right-most first.
	
	This is the combination of conventions we adopted (see pages~\pageref{def:function-application} 
	and~\pageref{def:fun_comp}). Outside of specialized areas of mathematics such 
	as the algebraic theory of automata or permutation group theory (or other branches of algebra as for example~\cite{Ferri2025_SplitLemmaGroupoids}), this seems
	to be the prevalent notion. For example, three famous textbooks from diverse fields: Lang's \emph{Algebra}~\cite{lang02},
	Spivak's \emph{Calculus on Manifolds}~\cite{spivak1971calculus} 
	or Königsberger's \emph{Analysis I}~\cite{koenigsberger2003analysis} (the latter being a standard text for first year math students in Germany, at least it was during the author's time) all use these conventions.
		
	\item Left-action and and left-most first.
		
	\label{prelim:discussion:Rrelation-fun-comp}
    The textbook by Almeida~\cite{Almeida95} from 1995 uses this convention,
	the survey article by Pin~\cite{DBLP:reference/hfl/Pin97} and, for example, a randomly sampled
	article by Diekert et al.~\cite{DBLP:journals/ijfcs/DiekertGK08}.
	Interestingly, function composition
	is not explicitly defined in these works.
	However, they have to use this convention, for otherwise, some cited results from this field
	have to be formulated differently. For example,
	partially ordered automata 
	%(this was first shown
	%\footnote{The used convention
		% 	is obvious from the proof of~\cite[Prop]{DBLP:journals/jcss/BrzozowskiF80}; we note that it contains
		% 	a typo, $\underline{x} = h$ should probably be $\underline{z} = h$. 
		%}
	%in~\cite{DBLP:journals/jcss/BrzozowskiF80})
	are precisely those automata
	having an $\mathcal R$-trivial transition monoid~\cite{DBLP:journals/jcss/BrzozowskiF80}  (we refer to the cited literature for definitions).
	But this result is only valid, due to the definition of the $\mathcal R$-relation via right ideals,
	when the left-most function is applied first in a composition (otherwise, it would correspond
	to $\mathcal L$-triviality).
	
	A recent (work-in-progress, available online) textbook by Pin~\cite{Pin2022} solved this by
	using two composition operations. To cite~\cite[Sect.~II.4.2]{Pin2022}: 
	% A recent (work-in-progress) textbook by Pin~\cite{} solved this by, basically,
	% mixing the conventions\footnote{Implicitly, we also mixed these conventions by letting words, via 
		% 	the function $\delta \colon Q \times \Sigma^* \to Q$, act on the right, whereas
		% 	elements of the transformation monoid, as functions, act on the left. But we 
		%    do not use two composition operations in the full transformation semigroup.\todo{monoid?}
		%}
	\begin{quote}
		[...] the reader should be warned that the product $\alpha\beta$ is not equal to $\alpha \circ\beta$,
		but to $\beta \circ \alpha$. In other word, the operation in $\mathcal T(P)$
		is reverse composition. [...]
	\end{quote}   
	
	% Pin is working on a textbook~\cite{} and we note that in its current version (available as a preprint online, 
	% so might be fixed
	% in the final version) it confuses both notions: in the beginning (via composition of relations)
	% the left-action/right-most first convention is adapted but later, when the mentioned result related $\mathcal R$-trivial
	% transition monoids is proven, by using words which act on the right (as is done here as well) 

	Another reason to adopt this convention
	is that it uses the common way of writing functions on the left, i.e., $f(x)$,
	while retaining compatibility with the \emph{composition of relations} 
	when defined in the following way: for a relation~$R$ from $X$ to $Y$ and a relation $S$ from $Y$ to $Z$,
	the composition is the relation $RS$ defined,
	for all $x \in X$ and $z \in Z$, by
	\[
	x(RS)z \Leftrightarrow \exists y \in Y \colon xRy \land ySz.
	\] 
		
	Now, when associating with $f \colon X \to Y$ the relation~$R_f$ from $X$ to $Y$
	by $xR_fy$ if and only if $\applyfunc{x}{f} = y$
	we have $R_{(f ; g)} = R_f R_g$
	for functions $f \colon X \to Y$, $g \colon Y \to Z$.	 
	This is, in particular, better suited when functions are defined 
	as relations (which we have not done in the present work since we introduced function
	before as the more primitive notion as in~\cite{Eilenberg1974}),
	or when representing relations by matrices, as this convention 
	fits with matrix multiplication.
	For example, in teaching first-year computer science students at Trier University~\cite{HenningDSSkript} 
	these conventions are adopted, as relations are introduced
	as a more primitive concept and (partial) functions are defined as special relations.
		
	Of course, one can define the composition of relations by reversing
	the order (this is done in~\cite[Sect.~I.1.2]{Pin2022}). However, in this case compatibility
	with matrix multiplication mentioned in the previous paragraph
	is lost (and, to the best of our knowledge, matrix multiplication is defined in the same way in all of mathematics)
	and a correct formula would involve additional matrix transpositions.
		
	\item Right-action and right-most first. We have never seen this combination in use.
\end{enumerate}

Furthermore, given an automaton $\semiautomat{A}{Q}{\Sigma}{\delta}$,
some authors choose to suppress 
mentioning~$\delta$ explicitly in the specification and application
of the transition function,
writing $q.u$ for $\applytrans{\delta}{q}{u}$ with $q \in Q$ and $u \in \Sigma^*$.
This dot-notation actually fits well with the right-action and left-most first conventions,
as then $q.u = (q) \fword{\delta}{u} = (q) (\fword{\delta}{v} ; \fword{\delta}{w}) = ((q) \fword{\delta}{v}) \fword{\delta}{w}$
for $u = vw$, $u,v,w \in \Sigma^*$, $q \in Q$.
However, we have opted against this dot-notation mainly because in \cref{part:constrsync}
often two transition functions (of a constraint automaton
and of an input automaton) appear in the same context. %concurrently, 
Hence we felt the need to distinguish them notation-wise.
Then, in our opinion, when applying
a transition function $\delta \colon Q \times \Sigma^* \to Q$ explicitly,
the left action convention seems to fit better.
Also, our convention adopted for composition was motivated by what
we feel is most often used in all of theoretical computer science and formal language theory (apart from specialized areas such as the algebraic theory of automata).

Note that ``operators'' such as $\Syn(\mathcal A)$ or $Z(G)$, which are never composed with other ``operators'', 
are almost always written as acting from the left, resulting in a mix of both conventions
in some works as for example in~\cite{DBLP:journals/tcs/AraujoBC13,Kurzweil04}.

Since \cref{part:aut} of the present thesis involves permutation group theory, we
were faced with the choice what convention to adopt. In face of \cref{part:constrsync}, which
does not involve any permutation group theory at all and is a more classical formal language theory topic,
we opted to the left-action convention that is more common in formal language theory.		
This convention has implications on how we define conjugation later. For more discussion on the two conventions of composition, 
see~\cite{nlab:composition} (the notation $f ; g$ is taken from this article).